\renewcommand{\Anno}{2010}

\renewcommand{\NumeroProblema}{1}

\renewcommand{\Testo}{
Siano $A$ e $B$ due punti distinti su una retta.

Determinare per quali interi positivi $n$ esistono $n$ punti distinti sulla 
retta, esterni rispetto al segmento $AB$, tali che la somma delle 
loro distanze da $A$ \`{e} uguale alla somma delle loro distanze da $B$.
}

\renewcommand{\Risposta}{
Esistono se e solo se $n$ \`{e} pari.
}

\renewcommand{\Aiutino}{
Per $n$ pari basta disporre opportunamente (come?) met\`{a} dei punti 
da una parte e met\`{a} dall'altra.

Per $n$ dispari consideriamo un punto qualunque $P$ esterno al
segmento.  Quali valori pu\`{o} assumere la differenza tra la sua
distanza da $A$ e la sua distanza da $B$? Cosa succede quando 
sommiamo questi valori su tutti i $P$?
}

\renewcommand{\SoluzioneUfficiale}{\TitoloSoluzione 
Gli $n$ punti con la propriet\`{a} richiesta esistono se e solo se 
$n$ \`{e} pari. Per dimostrarlo distinguiamo due casi.

\subparagraph{\textmd{Caso $n$ pari}}

Sia $n=2k$.  Scegliamo $k$ punti distinti a caso $P_{1},\ldots,P_{k}$
esterni rispetto al segmento $AB$ dalla parte di $A$.  Indichiamo con
$Q_{1},\ldots,Q_{k}$ i simmetrici di tali punti rispetto al punto
medio del segmento $AB$, come nella Figura~\ref{fig:puntini-PQ}.

\begin{figure}[htbp]
	\centering
	\psset{unit=4ex}
	\pspicture(-8,-1)(8,0.4)
	\psline(-8,0)(8,0)
	\pstGeonode[PointName=none,PointSymbol=|,dotscale=2](0,0){O}
	\pstGeonode[PosAngle=-90,PointSymbol=o,dotscale=2,linecolor=blue](-1,0){A}(1,0){B}
	\psset{dotscale=1.5,linecolor=red}
	\pstGeonode[PosAngle=-90,PointName=Q_{1}](2.2,0){Q1}
	\pstGeonode[PosAngle=-90,PointName=Q_{2}](7.5,0){Q2}
	\pstGeonode[PosAngle=-90,PointName=Q_{3}](3.5,0){Q3}
	\pstGeonode[PosAngle=-90,PointName=Q_{4}](5.6,0){Q4}
	\psset{dotscale=1.5,linecolor=olive}
	\pstRotation[PosAngle=-90,PointName=P_{1},RotAngle=180]{O}{Q1}[P1]
	\pstRotation[PosAngle=-90,PointName=P_{2},RotAngle=180]{O}{Q2}[P2]
	\pstRotation[PosAngle=-90,PointName=P_{3},RotAngle=180]{O}{Q3}[P3]
	\pstRotation[PosAngle=-90,PointName=P_{4},RotAngle=180]{O}{Q4}[P4]
	\endpspicture
	\caption{esempio di scelta dei $P_{i}$ e dei corrispondenti $Q_{i}$}
	\label{fig:puntini-PQ}
\end{figure}

Mostriamo che i $2k$ punti $P_{1},\ldots,P_{k},Q_{1},\ldots,Q_{k}$
soddisfano le propriet\`{a} richieste.  \`{E} intanto evidente che i
$Q_{1},\ldots,Q_{k}$ sono esterni rispetto al segmento $AB$ dalla
parte di $B$.  Inoltre, per ragioni di simmetria, per ogni
$i=1,\ldots,k$ si ha che
$$\mathrm{dist}(P_{i},A)=\mathrm{dist}(Q_{i},B) \quad\mbox{e}\quad
\mathrm{dist}(P_{i},B)=\mathrm{dist}(Q_{i},A),$$
e di conseguenza
$$\sum_{i=1}^{k}\mathrm{dist}(P_{i},A)+\sum_{i=1}^{k}\mathrm{dist}(Q_{i},A)=
\sum_{i=1}^{k}\mathrm{dist}(Q_{i},B)+\sum_{i=1}^{k}\mathrm{dist}(P_{i},B).$$

Abbiamo cos\`{i} mostrato che la somma delle distanze dei $2k$ punti
da $A$ \`{e} uguale alla somma delle distanze dei $2k$ punti da $B$,
il che completa la dimostrazione.

\subparagraph{\textmd{Caso $n$ dispari}}

Sia ora $n=2k+1$ un intero dispari.  Osserviamo che per ogni punto $P$
esterno rispetto al segmento $AB$ si ha che
$$\mathrm{dist}(P,A)-\mathrm{dist}(P,B)=\pm\mathrm{dist}(A,B),$$
dove si ha il segno $+$ od il segno $-$ a seconda che $P$ stia dalla 
parte di $B$ o dalla parte di $A$. Se ora abbiamo $2k+1$ punti 
$P_{1},\ldots,P_{2k+1}$, esterni rispetto al segmento $AB$, avremo 
che
$$\sum_{i=1}^{2k+1}\mathrm{dist}(P_{i},A)-
\sum_{i=1}^{2k+1}\mathrm{dist}(P_{i},B)=
\sum_{i=1}^{2k+1}\left(\mathrm{dist}(P_{i},A)-\mathrm{dist}(P_{i},B)\right).$$

Quella scritta al secondo membro \`{e} una somma con un numero 
dispari di addendi, ciascuno dei quali \`{e} uguale a 
$\pm\mathrm{dist}(A,B)$. \`{E} pertanto evidente che tale somma non 
potr\`{a} mai essere uguale a zero, e quindi le due sommatorie al 
primo membro non potranno mai essere uguali.
}

\renewcommand{\SoluzioniAlternative}{}

\renewcommand{\Osservazioni}{\TitoloOsservazione

Si noti che la condizione che i punti richiesti sono distinti \`{e}
irrilevante nella dimostrazione dell'impossibilit\`{a} per $n$
dispari.  }

\renewcommand{\Commenti}{
}